\section{The full Weil form on the actual theta input}
\label{sec:xi4-actual-theta-weil}

The compact spline calculation above does not test the noncompact input
whose Fourier transform is \(\Xi\). We now evaluate the complete Weil form
on that input after the actual multiplier \(w\). This section retains both
pole terms, the entire gamma integral, every prime power, all cross terms,
and explicit infinite-tail bounds. It uses the classical explicit formula
\eqref{eq:rh-full-weil-test}, with the conventions of Weil
\cite{weil1952} and Alp\"oge--Furman~\cite[Sections 1.1--1.2 and 2]{alpogeFurman2026}.

\subsection{The input, its involution, and the exact scale}

Put
\[
 \psi_0(x)=(4\pi^2x^4-6\pi x^2)e^{-\pi x^2},\qquad
 g_0(u)=e^{u/2}\sum_{n\ge1}\psi_0(ne^u).
\]
The Gaussian differentiation in \eqref{eq:xp-Hermite-seed} gives
\(\widehat\psi_0=\psi_0\) and \(\psi_0(0)=0\). Poisson summation therefore
gives \(g_0(-u)=g_0(u)\). Writing the source theta kernel as
\[
 \Phi(v)=\sum_{n\ge1}
 (2\pi^2n^4e^{9v}-3\pi n^2e^{5v})e^{-\pi n^2e^{4v}},
\]
one has, by direct substitution,
\begin{equation}\label{eq:xt-g0-scale}
 g_0(u)=\sum_{n\ge1}(4\pi^2n^4e^{9u/2}-6\pi n^2e^{5u/2})e^{-\pi n^2e^{2u}}
       =2\Phi(u/2).
\end{equation}
Since \(H_0(z)=\int_0^\infty\Phi(v)\cos(zv)\,dv\) and
\(\Xi(s)=8H_0(2s)\), the change of variable \(u=2v\) proves
\[
 \mathscr Fg_0(s)=\int_{\mathbb R}g_0(u)e^{-isu}\,du=\Xi(s).
\]
No rescaling of the source coefficient or of \(w\) is made.

Set \(q(u)=w(u)g_0(u)\) and \(Q=\mathscr Fq=\mathcal T_{\mathcal D}\Xi\).
The relation \(g_0(-u)=g_0(u)\) and \(w(-u)=\overline{w(u)}\) gives
\(q^\star(u):=\overline{q(-u)}=q(u)\). Hence
\[
 h(x):=(q*q^\star)(x)=\int_{\mathbb R}q(v)q(x-v)\,dv,
 \qquad \widehat h(z)=Q(z)^2.
\]
On the real axis \(Q(t)\) is real, so the Weil value is the real number
\begin{equation}\label{eq:xt-full-value}
 W(q,q)=2\Re Q(i/2)^2+\int_{\mathbb R}Q(t)^2G(t)\,dt
 -\sum_{n\ge2}\frac{\Lambda(n)}{\sqrt n}
       \{h(\log n)+h(-\log n)\},
\end{equation}
where
\[
 G(t)=\frac{1}{2\pi}\left(\Re\psi(1/4+it/2)-\log\pi\right).
\]
The correlation in \eqref{eq:xt-full-value} is not a pointwise value of
\(Q\); it is the full convolution and therefore retains every cross term.

\subsection{Uniform theta and correlation bounds}

For \(u\ge0\) write \(\Theta_N(u)\) for the first \(N\) terms of
\eqref{eq:xt-g0-scale}. Every summand is nonnegative, because
\(2\pi n^2e^{2u}-3\ge2\pi-3>0\). For \(N=8\), the omitted tail is bounded by
\begin{equation}\label{eq:xt-theta-tail}
 0\le g_0(u)-\Theta_8(u)\le
 E_8:=\frac{4\pi^2 9^4e^{-\pi9^2}}
 {1-(10/9)^4e^{-\pi(19)}}.
\end{equation}
Indeed \(e^{9u/2}e^{-\pi n^2e^{2u}}\le e^{-\pi n^2}\) for \(u\ge0\),
and the ratio of consecutive \(n^4\) majorants is at most the denominator
ratio displayed. The negative quadratic term can only reduce a summand,
so it need not be added to this upper bound. The full positive kernel has
the bound
\begin{equation}\label{eq:xt-theta-max}
 |g_0(u)|\le C_\Phi e^{-\pi},\quad
 C_\Phi=\frac{4\pi^2}{1-16e^{-3\pi}}
       +\frac{6\pi}{1-4e^{-3\pi}}.
\end{equation}

For \(x=\log n\ge0\), split the convolution integral into
\([-3,0]\), \([0,x]\), and \([x,x+3]\). On the first and last cells use the
evenness of \(g_0\) to replace negative arguments by positive ones; the
remaining two rays have one argument of size at least \(3\). Replacing
both factors by \eqref{eq:xt-theta-max} and integrating the elementary
majorants gives the exterior error
\begin{equation}\label{eq:xt-exterior-tail}
 E_{\rm ext}(x)=
 \frac{2C_\Phi^2 e^{-\pi}\exp(\tfrac92\,3-\pi e^6)
       (B+(x+3)/336)^2}
      {2\pi e^6-\tfrac92-
       [168(B+(x+3)/336)]^{-1}},\qquad B=|c_0|+c_1+c_2.
\end{equation}
Replacing \(g_0\) by \(\Theta_8\) on the three cells changes the product by
at most
\begin{equation}\label{eq:xt-trunc-tail}
 E_{\rm trunc}(x)=(x+6)(B+(x+3)/336)^2(2C_\Phi e^{-\pi}E_8+E_8^2).
\end{equation}
The bounds use only \(|w(u)|\le B+|u|/336\) and the product identity
\((a+\delta)(b+\delta)-ab=\delta(a+b)+\delta^2\).
The three finite cell integrals are evaluated by Arb interval
quadrature; every cell is split at all truncated-power knots. Thus the
reported ball plus \(E_{\rm ext}+E_{\rm trunc}\) encloses the exact \(h(x)\).

For the prime tail use a bound on the whole convolution line. Positivity
and the ratio of successive theta terms give, for every real \(u\),
\[
 |g_0(u)|\le C e^{9|u|/2-\pi e^{2|u|}},\qquad
 C=\frac{4\pi^2}{1-16e^{-3\pi}}<40.
\]
Put \(x=\log n\), \(v=x/2+y\), and \(A_n=B+x/672\).
The identities and inequalities
\[
\begin{gathered}
 |v|+|x-v|=\max(x,2|y|)\le x+2|y|,\\
 e^{2|v|}+e^{2|x-v|}
 \ge e^{2v}+e^{2(x-v)}=2n\cosh(2y)\ge2n+4ny^2,\\
 |w(v)w(x-v)|\le(A_n+|y|/336)^2,\qquad
 9|y|\le2\pi n y^2+\frac{81}{8\pi n}
\end{gathered}
\]
bound \(|h(x)|\) by
\[
C^2n^{9/2}e^{-2\pi n+81/(8\pi n)}
 \int_{\mathbb R}(A_n+|y|/336)^2e^{-2\pi n y^2}\,dy.
\]
Use \((a+b)^2\le2a^2+2b^2\) and the exact Gaussian moments
\[
 \int_{\mathbb R}e^{-2\pi n y^2}dy=(2n)^{-1/2},\qquad
 \int_{\mathbb R}y^2e^{-2\pi n y^2}dy
       =\frac1{4\pi n\sqrt{2n}}.
\]
The resulting coefficient multiplying \(A_n^2n^4e^{-2\pi n}\) is
\[
 C^2\sqrt2\,e^{81/(8\pi n)}
       \left(1+\frac1{4\pi n\,336^2A_n^2}\right).
\]
For \(n\ge17\), use \(A_n\ge\log n/672\), \(\log17>2\),
\(\pi>3\), and \(\sqrt2<10/7\).
Also \(81/(8\pi n)<1/5\) and \(e^{1/5}<5/4\), the latter obtained
by bounding its Taylor series by the geometric series.
Thus this coefficient is less than
\[
 1600\frac{10}{7}\frac54\frac{205}{204}
       =\frac{4100000}{1428}<3000.
\]
This proves the conservative envelope
\begin{equation}\label{eq:xt-prime-envelope}
 |h(\log n)|\le3000\left(B+\frac{\log n}{672}\right)^2n^4e^{-2\pi n}.
\end{equation}
The right side times \(\Lambda(n)/\sqrt n\) is decreasing for integers
\(n\ge17\): its consecutive ratio is at most
\[
 r=e^{-2\pi}(18/17)^{7/2}\frac{\log18}{\log17}
 \left(\frac{B+\log18/672}{B+\log17/672}\right)^2<1.
\]
Therefore the omitted prime-power sum after integer cutoff \(16\) is at most
\begin{equation}\label{eq:xt-prime-tail}
 E_{\rm prime}=6000\,17^{7/2}\log17
 \left(B+\frac{\log17}{672}\right)^2e^{-34\pi}/(1-r).
\end{equation}

For the gamma integral, shift the two theta Fourier contours by \(\pm i/2\).
The resulting exponential factor is \(a=\pi\cos1>0\); summing the geometric
theta majorants gives
\[
 C_\theta=\frac{4\pi^2}{1-16e^{-3a}}
           +\frac{6\pi}{1-4e^{-3a}}.
\]
The polynomial factor and one derivative term in \(w\) give the explicit
constant
\[
 C_K=2C_\theta e^{13/4}
 \left(\frac{9}{2ae}\right)^{9/4}e^{-a/2}
 \left(\frac{B+1/672}{a}+\frac1{336a^2}\right).
\]
The standard digamma bound
\[
 |G(t)|\le\{\gamma_E+8+\log\pi+\log(1+2|t|)\}/(2\pi)
\]
then yields, for \(T=64\),
\begin{equation}\label{eq:xt-gamma-tail}
 E_\Gamma(T)=\frac{C_K^2}{\pi}e^{-T}
 \left(\gamma_E+8+\log\pi+\log(1+2T)+\frac1{T+1/2}\right).
\end{equation}
Here are the contour details. For \(u\ge0\), the theta series on
\(u\pm i/2\) has absolute value at most
\(C_\theta e^{9u/2-ae^{2u}}\); evenness supplies the other half-line.
The multiplier on that strip is at most
\(e^{13/4}(B+1/672+|u|/336)\).
Split the theta exponential into two equal factors and use
\[
 e^{9u/2-(a/2)e^{2u}}\le
       \left(\frac9{2ae}\right)^{9/4},\qquad
 e^{-(a/2)e^{2u}}\le e^{-a/2-au}.
\]
Integrating \(e^{-au}\) and \(ue^{-au}\) gives exactly \(C_K\).
Horizontal contour displacement, with the vertical ends tending to zero
by these estimates, proves \(|Q(t)|\le C_K e^{-|t|/2}\).
For the digamma bound, use its convergent series at \(z=1/4+ib\):
\[
 \frac1n-\frac{n+1/4}{(n+1/4)^2+b^2}
 =\frac{1/4}{n(n+1/4)}
  +\frac{b^2}{(n+1/4)((n+1/4)^2+b^2)}.
\]
Splitting the second sum at \(\max(1,\lceil |b|\rceil)\) bounds its
tail by \(1/2\) and its initial part by \(1+\log(1+|b|)\).
The \(n=0\) real term has absolute value at most \(4\).
Together these imply the stated, larger bound for \(|G(t)|\).
Finally \(\log(1+2t)\le\log(1+2T)+(t-T)/(T+1/2)\) for \(t\ge T\).
Integrating \(e^{-t}\) and \((t-T)e^{-t}\) proves the displayed tail.
No zero ordinates or RH input occurs.

\subsection{The certified value}

The one-worker Arb calculation uses 192-bit balls, \(N=8\), \(T=64\), prime
integer cutoff \(16\), margin \(3\), and absolute/relative tolerance \(2^{-90}\).
It integrates all 128 gamma cells and the three correlation cells for each
prime power \(n\le16\), then adds \eqref{eq:xt-prime-tail},
\eqref{eq:xt-gamma-tail}, \eqref{eq:xt-exterior-tail} and
\eqref{eq:xt-trunc-tail}. The exact receipt is
\begin{center}\small\ttfamily
agents/connes\_actual\_input\_review/actual\_theta\_weil/actual\_theta\_weil\_results.json
\end{center}
with checker hash
\begin{center}\small\ttfamily
d0bbeb4cc91d0ff522655f8e8bf39f0d2981a48c41834e595dc5ffbd6c532279
\end{center}
The component enclosures are
\[
 \begin{aligned}
 \text{pole terms}&=[2.1389982672320680\,10^{-7}\pm4.11\,10^{-62}],\\
 \text{prime terms through }16&=[1.4681655033429546\,10^{-8}\pm4.62\,10^{-28}],\\
 E_{\rm prime}&<1.053\,10^{-43},\qquad E_\Gamma(64)<8.103\,10^{-27},
 \end{aligned}
\]
and the complete outward-rounded value is
\begin{equation}\label{eq:xt-result}
 W(wg_0,wg_0)=[2.0984855607004\,10^{-11}\pm6.08\,10^{-25}]>0.
\end{equation}
Thus the actual theta input is not a negative Weil witness. This positive
value is separate from the already proved negative linear coefficient
\(K(\gamma)\) and from the uncanceled pole \(K/\Xi\): one is a quadratic
full-form value and the other is a linear shifted evaluation. The
calculation establishes positivity for this one actual noncompact input,
not global Weil positivity and not RH.
